% =====================================================================
%  CARNET D'ENTRAÎNEMENT — FICHE 8 — THÈME 1 — NIVEAU FACILE — ÉNONCÉS
% =====================================================================
\documentclass[11pt,a4paper]{article}
\usepackage[utf8]{inputenc}
\usepackage[T1]{fontenc}
\usepackage{lmodern}
\usepackage[french]{babel}
\usepackage{amsmath,amssymb,mathtools}
\usepackage{icomma}
\usepackage{siunitx}
\sisetup{output-decimal-marker={,},per-mode=power,group-separator={\,},group-minimum-digits=5}
\usepackage[version=4]{mhchem}
\usepackage{xcolor}
\usepackage{colortbl}
\usepackage{tikz}
\usepackage{pgfplots}
\usepackage[most]{tcolorbox}
\usepackage{enumitem}
\usepackage{array}
\usepackage{booktabs}
\usepackage{multicol}
\usepackage{fancyhdr}
\usepackage[hidelinks]{hyperref}
\usetikzlibrary{arrows.meta,positioning,calc,shapes.geometric,fit}
\pgfplotsset{compat=1.18}
\usepackage[a4paper,margin=2.1cm,headheight=26pt,headsep=10pt]{geometry}
\usepackage{titlesec}

\definecolor{primary}{RGB}{150,98,162}
\definecolor{primarydark}{RGB}{92,52,110}
\definecolor{excol}{RGB}{0,135,90}
\definecolor{exocol}{RGB}{200,60,40}
\definecolor{endocol}{RGB}{0,110,180}

\tcbset{boxbase/.style={enhanced,breakable,boxrule=0.9pt,arc=2.5pt,left=3mm,right=3mm,top=2.3mm,bottom=2.3mm,fonttitle=\bfseries,coltitle=white,toptitle=0.6mm,bottomtitle=0.6mm}}
\newtcolorbox[auto counter]{exercice}[1][]{boxbase,colback=primary!4,colframe=primary,title={\textsc{Exercice}~\thetcbcounter\ \ifx\relax#1\relax\relax\else\textmd{--- \itshape #1}\fi}}

\newcommand{\dH}{\Delta H}
\newcommand{\dU}{\Delta U}

\pagestyle{fancy}
\fancyhf{}
\fancyhead[L]{\small\textcolor{primarydark}{\textbf{Fiche 8} — Thème 1 : Premier principe et enthalpie}}
\fancyhead[R]{\small\textcolor{primarydark}{\textsc{Facile}}}
\fancyfoot[C]{\small\thepage}
\renewcommand{\headrulewidth}{0.8pt}
\renewcommand{\headrule}{\hbox to\headwidth{\color{primary}\leaders\hrule height \headrulewidth\hfill}}
\setlength{\parindent}{0pt}\setlength{\parskip}{3pt}

\begin{document}

\begin{center}
\begin{tikzpicture}
\node[rectangle,rounded corners=7pt,fill=primary,text=white,minimum width=16cm,
      minimum height=1.1cm,align=center,inner sep=6pt]
  {{\large\bfseries Thème 1 — Premier principe et enthalpie}\\[1pt]{\small Niveau \textsc{Facile} — énoncés}};
\end{tikzpicture}
\end{center}

\begin{exercice}[conventions de signe]
Pour chaque situation, du \textbf{point de vue du système}, indiquer le signe (positif ou négatif) de la
grandeur demandée, en justifiant en une phrase.
\begin{enumerate}[label=\alph*),leftmargin=2em,topsep=2pt,itemsep=2pt]
  \item Le système \textbf{reçoit} \SI{300}{\joule} de chaleur du milieu extérieur. Signe de $Q$ ?
  \item Un gaz se \textbf{dilate} et pousse l'atmosphère en fournissant \SI{120}{\joule} de travail. Signe de $W$ ?
  \item Une réaction \textbf{exothermique} se déroule à pression constante. Signe de $\dH$ ?
  \item Le système est \textbf{isolé}. Que valent $Q$, $W$, puis $\dU$ ?
\end{enumerate}
\end{exercice}

\begin{exercice}[appliquer $\dU = Q + W$]
Un système fermé subit une transformation. Dans chaque cas, écrire d'abord $Q$ et $W$ \emph{avec leur signe},
puis calculer $\dU$ en joules.
\begin{enumerate}[label=\alph*),leftmargin=2em,topsep=2pt,itemsep=2pt]
  \item Il \textbf{reçoit} \SI{500}{\joule} de chaleur et \textbf{reçoit} \SI{150}{\joule} de travail.
  \item Il \textbf{reçoit} \SI{400}{\joule} de chaleur et \textbf{cède} \SI{250}{\joule} de travail.
  \item Il \textbf{cède} \SI{600}{\joule} de chaleur et \textbf{reçoit} \SI{600}{\joule} de travail.
\end{enumerate}
\end{exercice}

\begin{exercice}[exo ou endo à partir du signe de $\dH$]
Pour chacune des transformations suivantes, dont on donne la variation d'enthalpie, indiquer si elle est
\textbf{exothermique} ou \textbf{endothermique}, et préciser si les produits ont une enthalpie
\emph{plus haute} ou \emph{plus basse} que celle des réactifs.
\begin{enumerate}[label=\alph*),leftmargin=2em,topsep=2pt,itemsep=2pt]
  \item $\ce{N2}(g) + 3\,\ce{H2}(g) \rightarrow 2\,\ce{NH3}(g)$, \quad $\dH = \SI{-92}{\kilo\joule}$.
  \item $2\,\ce{H2O}(l) \rightarrow 2\,\ce{H2}(g) + \ce{O2}(g)$, \quad $\dH = \SI{+572}{\kilo\joule}$.
  \item $\ce{CH4}(g) + 2\,\ce{O2}(g) \rightarrow \ce{CO2}(g) + 2\,\ce{H2O}(l)$, \quad $\dH = \SI{-890}{\kilo\joule}$.
  \item $\ce{CaCO3}(s) \rightarrow \ce{CaO}(s) + \ce{CO2}(g)$, \quad $\dH = \SI{+178}{\kilo\joule}$.
\end{enumerate}
\end{exercice}

\begin{exercice}[lire un diagramme enthalpique]
On donne ci-dessous le diagramme enthalpique d'une réaction $\ce{A} \rightarrow \ce{B}$, avec les
enthalpies exprimées en \si{\kilo\joule}.
\begin{center}
\begin{tikzpicture}[scale=1.0]
  \draw[->,thick] (0,-0.2) -- (0,4.4) node[above,font=\small]{$H$ (kJ)};
  \draw[thick] (0,0) -- (7.6,0) node[right,font=\small]{avancement};
  % Réactifs à 120
  \draw[exocol,very thick] (0.6,3.6) -- (2.7,3.6);
  \node[anchor=south west,font=\small] at (0.65,3.65){Réactifs (A)};
  \node[anchor=east,font=\scriptsize] at (0.55,3.6){$120$};
  \draw[dashed,gray] (2.7,3.6) -- (4.2,3.6);
  % Produits à 40
  \draw[exocol,very thick] (4.2,1.2) -- (6.9,1.2);
  \node[anchor=south west,font=\small] at (4.25,1.25){Produits (B)};
  \node[anchor=east,font=\scriptsize] at (0.55,1.2){$40$};
  \draw[dashed,gray] (0,1.2) -- (4.2,1.2);
  % flèche
  \draw[->,exocol,line width=1.2pt] (3.6,3.6) -- (3.6,1.2) node[midway,right,font=\small]{$\dH$};
\end{tikzpicture}
\end{center}
\begin{enumerate}[label=\alph*),leftmargin=2em,topsep=2pt,itemsep=2pt]
  \item Calculer $\dH$ de la réaction (donner sa valeur et son signe).
  \item La réaction est-elle exothermique ou endothermique ?
  \item La chaleur est-elle \emph{dégagée} vers l'extérieur ou \emph{absorbée} ?
\end{enumerate}
\end{exercice}

\begin{exercice}[vocabulaire de base]
Répondre brièvement.
\begin{enumerate}[label=\alph*),leftmargin=2em,topsep=2pt,itemsep=2pt]
  \item Qu'est-ce qu'un \textbf{système isolé} ? Que valent alors $Q$, $W$ et $\dU$ ?
  \item Quelle est la différence entre une transformation \textbf{athermique} et une enceinte \textbf{adiabatique} ?
  \item À quel type de système (ouvert, fermé, isolé) un \textbf{calorimètre parfait} correspond-il ?
\end{enumerate}
\end{exercice}

\end{document}
